\documentclass[11pt]{scrartcl}
\usepackage[sexy]{evan}
\usepackage{answers}

\begin{document}
\title{Functional Equations Solutions}
\author{Prakrit Gajurel}
\date{9 July 2026}
\maketitle

\section{Walkthrough Problems}

\begin{problem}
    [USAMO 2002]
    Determine all functions $f: \mathbb{R} \rightarrow \mathbb{R}$ such that
    \[
    f(x^2-y^2) = xf(x) - yf(y)
    \]
    for all pairs of real numbers $x$ and $y$.
    
\end{problem}

\begin{proof}
    Let $P(x, y)$ denote the given assertion. Note that $P(0,0)$ gives us $f(0)=0$. Then,
    \begin{align*}
        P(x, -x) \Longrightarrow &f(x^2-x^2) = xf(x) + xf(-x)
        \\& 0 = x(f(x)+f(-x))
    \end{align*}
    For nonzero $x$, we have $f(-x) = -f(x)$, but when $x=0$ we have $f(0)=0$ anyway, so we have $f$ is odd for all real numbers $x$.
    \[
    P(x,0) \Longrightarrow f(x^2) = xf(x)
    \]
    So now we can rewrite the original condition as:
    \[
    f(x^2 - y^2) = f(x^2) - f(y^2)
    \]
    We let $a = x^2$, and $b = y^2$ such that, for nonnegative $a$, $b$ we have:
    \[
    f(a-b) = f(a) - f(b)
    \]
    Again, we can set $a=x_0 + y_0$ and $b=y_0$ to get:
    \[
    f(x_0+y_0-y_0) = f(x_0+y_0) - f(y_0) \implies f(x_0 + y_0) = f(x_0) + f(y_0)
    \]
    for nonnegative $x_0$, $y_0$. However, since $f$ is odd, we get that $f$ is additive on the negatives too.

    Finally, we let $x=a+b$ in one of the previous equations $f(x^2) = xf(x)$ to get:
    \[
    f(a^2 + 2ab + b^2) = (a+b)(f(a) + f(b))
    \]
    \[
      f(2ab) = af(b) + bf(a)
    \]
    \[
      a = 1 \implies f(2b) = f(b) + f(1)b = 2f(b)
    \]
    Which gives $f(x) = cx$ for all $x$, which upon checking works.
\end{proof}
\newpage
\begin{problem}
  [IMO 2017]
  Solve over $\mathbb{R}$ the functional equation
  \[
  f(f(x)f(y)) + f(x+y) = f(xy).
  \]
\end{problem}

\begin{proof}
  We claim that $f(x) = 0$, $f(x) = 1-x$, and $f(x) = x-1$ are the only solutions. It's easy to check that these solutions indeed work.
  \\ \\
  Let $P(x,y)$ denote the given assertion. First, note:
  \[
  P(0,0) \Longrightarrow f(f(0)^2 = 0) \implies \text{ there exists a } z \text{ such that } f(z) = 0.
  \]
  Also, $P(x, 0)$ gives that $f(0)=0 \implies f \equiv 0$. So from now assume that $f(0) \neq 0$. We take, for $z\neq1$,
  \[
  P\left(z, \frac{z}{z-1}\right) \Longrightarrow f\left(f(z)f\left(\frac{z}{z-1}\right)\right) = 0
  \]
  giving $f(0)=0$, which is a contradiction, giving $z=f(0)^2=1 \implies f(0) = \pm 1$. One can now take
  $P(1,x)$ to give $f(x+1) = f(x)-1$, then $P(x,0)$ to get $f(f(x)) = 1-f(x)$. Applying $f$ on this equation and equating
  some terms gives us:
  \[
  f(f(f(x))) = f(1-f(x)) = 1-f(f(x)) = f(x)
  \]
  Before we move on, I'd like to prove something that would help us in our final claim.
  \begin{claim}
    For sufficiently large $N_1$ and $N_2$, we can find $x$, $y$ satisfying $x+y = N_1 + 1$ and $xy = N_2$.
  \end{claim}
  \begin{proof}
    First we observe that we can think of $x$, $y$ as roots of some quadratic equation. Of course, comparing it with Vieta's gives us that
    $x$ and $y$ are possible roots of the equation
    \[
    x_0^2 - (N_1 + 1)x_0 + N_2 = 0
    \]
    Obviously distinct solutions $x$, $y$ to this equation exist if and only if the discriminant, $\Delta$ is greater than zero. So we have:
    \begin{align*}
    \Delta &= (N_1+1)^2 - 4N_2 \\
    &= N_1^2 + 2N_1 + 1 - 4N_2 > 0
    \end{align*}
    So we can just pick sufficiently large $N_1$, $N_2$ that satisfies the condition(or maybe something like $N_1^2 \ge 4N_2$) so we have 
    solutions for $x, y$.
  \end{proof}

  Now, we prove the final claim.

  \begin{claim}
    $f$ is injective.
  \end{claim}
  \begin{proof}
    Assume for some sufficiently large $a$, $b$, we have $f(a) = f(b)$ and also, $x+y = a+1$ and $xy = b$(from the above claim, we can do this).
    Then, the original assertion gives us
    \begin{align*}
      f(f(x)f(y)) &= f(b) - f(a+1) \\
      &= f(b) - (f(a) - 1) \\
      &= 1
    \end{align*}
    Shifting $f$ we get that $f(f(x)f(y)+1) = 0 \implies f(x)f(y) + 1 = 1$ so $f(x)f(y) = 0$, giving that $1 \in \left\{ x, y \right\}$. However,
    we are done as if $x=1$ we get $y+1=a+1 \implies y=a$ and $xy=b \implies b=y=a$. This follows similarly for $y=1$, and so $f$ is injective.
  \end{proof}

  From the injectivity, and our earlier calculations we get $f(1-f(x)) = f(x) \implies f(x) = 1-x$ and since $-f$ must be a solution too, 
  $f(x) = x-1$ is another solution. Therefore, our solutions are:

  \[
  \boxed{f(x) = 0, 1-x, x-1}
  \]
\end{proof}

\begin{problem}
  [USAMO 2018]
  Find all functions $f: (0, \infty) \rightarrow (0, \infty)$ such that
  \[
  f\left( x + \frac 1y \right) + f\left( y + \frac 1z \right) + f\left( z + \frac 1x \right) = 1
  \]
  for all $x$, $y$, $z > 0$ with $xyz = 1$.
\end{problem}

\begin{proof}
  We begin by substituting the values $x=\frac ab$, $y=\frac bc$, $z=\frac ca$ to get the following:
  \[
  \sum_{\text{cyc}} f\left( \frac{a+c}{b} \right) = 1
  \]
  We then provide another substitution, this time for the function. Define a function $g: (0, 1) \rightarrow (0, 1)$ such that
  $g(x) = f(\frac{1-x}{x})$, then for $a + b + c = 1$ it follows that
  \[
  g(a) + g(b) + g(c) = 1
  \]
  \begin{claim}
    $g$ satisfies Jensen's functional equation over $(0, \frac 12)$.
    \begin{proof}
      Define variables $x$, $y \in (0, \frac 12)$ such that $\frac{x+y}{2} \in (0, \frac 12) \subset (0, 1)$ and also $0 < x+y < 1$.
      Therefore, we can make the following substitutions:
      \[
      a = b = \frac{x+y}{2}
      \]
      \[
      c = 1-(x+y)
      \]
      Which satisfy $a+b+c=1$, then,
      \[
      g\left( \frac{x+y}{2} \right) + g\left( \frac{x+y}{2} \right) + g(1-(x+y)) = 1
      \]
      and
      \[
      g\left( x \right) + g\left( y \right) + g(1-(x+y)) = 1
      \]
      Give us that
      \[
      g(x) + g(y) = 2g\left( \frac{x+y}{2} \right),
      \]
      as desired.
    \end{proof}
  \end{claim}
  We now define a new function $h: [0, 1] \rightarrow \mathbb{R}$ as:
  \[
  h(t) = g \left( \frac{2t+1}{8} \right) - (1-t) g\left( \frac 18\right) - tg\left( \frac 38 \right)
  \]
  Notice that $h(0) = h(1) = 0$. Also,
  \begin{align*}
    h\left( \frac 12 \right) &= g \left( \frac 28 \right) - \frac 12 g\left( \frac 18 \right) - \frac 12 \left( \frac 38 \right) \\
    &= g \left( \frac 18 \right) + g \left( \frac 18 \right) - \frac 12 g \left( \frac 18 \right) - \frac 12 g \left( \frac 18 \right) - \frac 12 g \left( \frac 18 \right) \\
    &= 0
  \end{align*}
  So now we have $h(0) = h(1) = h(\frac 12) = 0$.
\end{proof}

\newpage{}

\section{Problems}

\begin{problem}
  [IMO 2008]
  Find all functions $f$ from the positive reals to the positive reals such that
  \[
  \frac{f(w)^2 + f(x)^2}{f(y^2) + f(z^2)} = \frac{w^2 + x^2}{y^2 + z^2}
  \]
  for all positive real numbers $w$, $x$, $y$, $z$ satisfying $wx=yz$.
\end{problem}

\begin{proof}
  We let $P(w, x, y, z)$ represent the given assertion. We claim that $f(x) = x$ and $f(x) = \frac 1x$ are the only solutions, it is
  easy to check that these work. \newline \newline
  First, note
  \[
  P(x, x, x, x) \Longrightarrow \frac{f(x)^2}{f(x^2)} = 1 \implies f(x^2) = f(x)^2
  \]
  Which also gives $f(1) = 1$. Further,
  \[
  P(x, 1, \sqrt x, \sqrt x) \Longrightarrow (f(x)-x)(xf(x)-1) = 0
  \]
  Which gives us that $f(x) = x$ or $\frac 1x$. However, as $f$ could be a piecewise function, we prove the nonexistence of one.o
  \newline \newline
  We assume $f(a) = a$ and $f(b) = \frac 1b$ for some distinct $a$, $b$. We note:
  \[
  P(\sqrt a, \sqrt b, \sqrt{ab}, 1) \Longrightarrow \frac{a+\frac 1b}{f(ab) + 1} = \frac{a+b}{ab + 1}
  \]
  If we take $f(ab) = ab$, we get $b=1$ and the case where $f(ab) = \frac{1}{ab}$ goves $a=1$, contradicting that $a$ and $b$ are distinct.
  Thus, $f$ is not piecewise, and the only solutions are $f(x) = x$ and $f(x) = \frac 1x$ for all $x \in \mathbb{R}^+$.
\end{proof}

\begin{problem}
  [IMO 2010]
  Find all functions $f: \mathbb{R} \rightarrow \mathbb{R}$ such that for all $x$, $y \in \mathbb{R}$,
  \[
  f(\lfloor x \rfloor y) = f(x) \lfloor f(y) \rfloor
  \]
\end{problem}

\begin{proof}
  Let $P(x, y)$ denote the given assertion. Note that
  \[
  P(0, 0) \Longrightarrow f(0) = f(0) \lfloor f(0) \rfloor
  \]
  From which we get that either $f(0) = 0$, or $1 \le f(0) < 2$, more specifically we get $\lfloor f(0) \rfloor = 1$. We now
  consider cases.
  \newline \newline
  $\boxed{\textbf{CASE 1.}}$ $f(0) = 0$
  Let some $x$ be such that $0 < x < 1$, then the original equation becomes:
  \[
  f(x) \lfloor f(y) \rfloor = 0
  \]
  Which gives us either $\lfloor f(y) \rfloor = 0$ or $f(x) = 0$ for all $0<x<1$. For the first case, consider now
  $P(1, y)$ giving us $f(y) = 0$ for all $y$.
\end{proof}\

\begin{problem}
  [IMO 2009]
  Find all functions $f: \mathbb{Z}_{>0} \rightarrow \mathbb{Z}_{>0}$ such that for positive integers $a$ and $b$, the numbers
  \[
  a, \qquad f(b), \qquad, f(b+f(a)-1)
  \]
  are the sides of a non-degenerate triangle.
\end{problem}

\begin{proof}
  We prove something that will be crucial.
  \begin{claim}
    If $1$, $a$, $b$ are sides of a triangle where $a$, $b \in \mathbb{Z}_{>0}$, then $a=b$.
    \begin{proof}
      By the triangle inequality,
      \[
      b < a + 1
      \]
      \[
      a < 1 + b
      \]
      Combining gives us $-1 < a-b < 1$. Since $a$, $b$ are positive integers, we get that $a-b=0 \implies a=b$.
    \end{proof}
  \end{claim}
  Let $P(a, b)$ denote the tuple $(a, f(b), f(b+f(a)-1))$. Then,
  \[
  P(1, b) \Longrightarrow (1, f(b), f(b+f(1)-1))
  \]
  From our above claim, we get that $f(b) = f(b+f(1)-1)$, so $f(1) \neq 1 \implies $$f$ is periodic, and takes on finitely many values. 
  Note that this also means $f$ has a maximum(again, if and only if $f(1) \neq 1$).
  \begin{claim}
    $f(1) = 1$.
    \begin{proof}
      We proceed by contradiction, assume not, then obviously $f(1) > 1$. We choose a sufficiently large $a$ for
      $P(a, b)$, then by the triangle inequality,
      \[
      a < f(b+f(a)-1) - f(b)
      \]
      Which would hold false.
    \end{proof}
  \end{claim}
\end{proof}

\end{document}